% ===== Wrap-up (~4 min, then questions) =====
\partaccent{peach}\notopic

\begin{frame}{Այսօրվա 5 գաղափարը}
  \begin{enumerate}
    \item<1-> IQ-ն խելքի «քանակը» չէ, այլ քո \textbf{տեղը մյուսների համեմատ}. միջինը միշտ 100 է
    \item<2-> Կա ընդհանուր խելք (g), ու կա \textbf{հեղուկ} (նոր խնդիրներ) և \textbf{բյուրեղացած} (գիտելիք) ինտելեկտ
    \item<3-> Քիչ գիտեցողները հաճախ ավելի \textbf{վստահ} են, խելացիները կասկածում են, իսկ մենք լսում ենք վստահներին
    \item<4-> Խելացի ուղեղը ոչ թե ավելի շատ է աշխատում, այլ ունի \textbf{արդյունավետ կապեր}։ «Ուղեղի մարզման» խաղերը լավացնում են միայն իրենց
    \item<5-> Խելքը մասամբ ժառանգական է, բայց կախված է \textbf{մշակույթից, սպասումներից ու կրթությունից}
  \end{enumerate}
\end{frame}

\begin{frame}{Փորձիր տանը}
  \begin{columns}[T]
    \begin{column}{0.31\textwidth}
      \begin{card}[colback=mint!30!cardwhite]
        \centering\bigemoji[1cm]{memo}\par\smallskip
        \raggedright Հաջորդ քննությունից առաջ գրի առ, թե ինչ գնահատական ես սպասում։ Հետո համեմատիր։ Այսպես ես մարզում \textbf{ինքնագնահատումը}
      \end{card}
    \end{column}
    \begin{column}{0.31\textwidth}
      \begin{card}[colback=lavender!35!cardwhite]
        \centering\bigemoji[1cm]{guitar}\par\smallskip
        \raggedright Սովորիր իսկապես \textbf{նոր} բան՝ գործիք, լեզու, խաղ։ Հազարերորդ խաչբառը չի օգնի
      \end{card}
    \end{column}
    \begin{column}{0.31\textwidth}
      \begin{card}[colback=peach!35!cardwhite]
        \centering\bigemoji[1cm]{speakinghead}\par\smallskip
        \raggedright Հաջորդ վեճում նկատիր, թե ով է ամենավստահը։ Նա նաև ամենաճի՞շտն է
      \end{card}
    \end{column}
  \end{columns}
\end{frame}

\begin{frame}{Հաջորդ վեբինարը}
  \begin{columns}[c]
    \begin{column}{0.6\textwidth}
      \chip[lavender]{eyes}{5. Տեսա՞ր, որ այս գլուխը գալիս է}\par\smallskip
      Ինչպես է ուղեղը ընկալում աշխարհը.\par\smallskip
      \begin{itemize}
        \item \emoji{nose}\ ինչու է հոտը համից ուժեղ
        \item \emoji{ear}\ ինչպես են լսողությունն ու շոշափումը կապված
        \item \emoji{bread}\ ինչու ենք դեմքեր տեսնում տոստի վրա
        \item \emoji{speech}\ ինչու չենք կարողանում ականջ չդնել ուրիշների խոսակցությանը
      \end{itemize}
      \medskip
      {\small\color{plum!70!cream}Այնուհետև. 6. Անհատականություն, 7. Ուրիշ մարդիկ, 8. Երբ ուղեղը «փչանում» է}
    \end{column}
    \begin{column}{0.34\textwidth}
      \centering\mascot[2.8cm]{confused}
    \end{column}
  \end{columns}
\end{frame}

\begin{frame}[plain]
  \begin{tikzpicture}[remember picture, overlay]
    \fill[lavender!60!cream] (current page.south west) rectangle (current page.north east);
    \fill[pattern={Dots[distance=4.4mm, radius=0.28mm]}, pattern color=plum!10!lavender!60!cream]
      (current page.south west) rectangle (current page.north east);
    \node at ([xshift=-3.2cm]current page.center) {\mascot[4.4cm]{happy}};
    \node[anchor=west, text=plum] at ([xshift=0.2cm, yshift=0.6cm]current page.center)
      {\titlefont\fontsize{40}{42}\selectfont Հարցեր};
    \node[anchor=west, text=plum, font=\mediumfont\large, align=left] at ([xshift=0.2cm, yshift=-1.0cm]current page.center)
      {Գրիր չաթում կամ\\ միացրու խոսափողը \emoji{megaphone}};
    \node[anchor=west, text=plum!70!cream, font=\small] at ([xshift=0.2cm, yshift=-2.1cm]current page.center)
      {Դին Բըրնեթ, «The Idiot Brain», 2016};
  \end{tikzpicture}
\end{frame}
